% ============================================================
% Literature Review
% ============================================================

\section{Conceptual Framing}

This chapter reviews theories and empirical studies that explain how colour design in floral gift contexts can influence users' emotions and downstream intentions. The review is organized around three layers: (1) colour as a perceptual and compositional design factor, (2) emotion as the core psychological mechanism, and (3) behavioral intentions as outcome variables. Finally, this chapter translates prior work into a concrete research design that combines generative image conditions with quantitative and qualitative measures.

\section{Colour}

\subsection{Colour as a Focal Design Factor}

Colour is one of the earliest and most salient visual cues processed by viewers. In product and interface evaluation, colour can rapidly shape first impressions, perceived quality, and affective tone before semantic interpretation is complete. In floral gift design, colour is especially central because flowers are already culturally encoded through hue-symbol associations (e.g., romance, celebration, sympathy), making colour choices both aesthetic and communicative.

\subsection{Primary--Secondary--Accent Structure}

Design studies commonly distinguish a layered colour composition strategy consisting of primary, secondary, and accent colours. The primary colour establishes the dominant atmosphere, the secondary colour supports balance and depth, and the accent colour guides attention to focal elements. This structure helps avoid monotony while preserving coherence, and it provides an operational basis for manipulating colour configurations in controlled experiments.

\subsection{Colour Space and Color Attributes}

To support rigorous analysis, colour should be represented in a formal colour space (e.g., RGB, HSV/HSL, CIELAB). In this study, we adopt the CIE $L^\ast a^\ast b^\ast$ (CIELAB) framework because it keeps metric comparability while making perceptual interpretation clearer for floral-material colour analysis.

Commonly used colour representation methods in design research include color models, color spaces, and practical color-order systems. Compared with directly using raw RGB coordinates, CIELAB better separates perceptual dimensions and supports quantitative comparisons across generated designs.

\begin{table}[htbp]
  \centering
  \begin{threeparttable}
    \caption{Color attributes and derived metrics.}
    \label{tab:color_metrics}
    \renewcommand{\arraystretch}{1.25}
    \begin{tabularx}{0.96\textwidth}{>{\raggedright\arraybackslash}p{0.16\textwidth}>{\raggedright\arraybackslash}p{0.14\textwidth}>{\raggedright\arraybackslash}p{0.30\textwidth}>{\raggedright\arraybackslash}X}
      \toprule
      \textbf{Color Attribute} & \textbf{Symbol} & \textbf{Definition (Brief)} & \textbf{Analytical Use in This Study} \\
      \midrule
      Lightness & $L^\ast$ & Perceived light--dark dimension in CIELAB & Quantifies tonal layering and lightness contrast across materials and the whole composition. \\
      Chroma & $C^\ast_{ab}$ & Colorfulness strength derived from $a^\ast, b^\ast$ & Captures saturation-like intensity and chroma control of floral palettes. \\
      Hue angle & $h_{ab}$ & Hue direction on a $0^\circ$--$360^\circ$ circle in $a^\ast$--$b^\ast$ plane & Describes hue organization (dominant hue) and hue differences among design elements. \\
      Color difference & $\Delta E_{00}$ & Perceptual distance between two colors & Measures inter-element and inter-design color distances for comparability checks. \\
      Contrast (derived) & Contrast & Composite contrast from $L^\ast$, $C^\ast_{ab}$, $h_{ab}$ and/or $\Delta E_{00}$ & Quantifies visual separation and salience in generated outcomes. \\
      Harmony (derived) & Harmony & Composite harmony metric based on hue/lightness/chroma relations & Quantifies perceived coherence and compatibility of multi-color compositions. \\
      \bottomrule
    \end{tabularx}
    \begin{tablenotes}[flushleft]
      \footnotesize
      \item Definitions follow the CIE 1976 CIELAB framework; color-difference values use CIEDE2000 where applicable.
    \end{tablenotes}
  \end{threeparttable}
\end{table}

As shown in \autoref{tab:color_metrics}, combining base attributes and derived metrics enables reproducible color characterization and supports linking computational descriptors with human evaluation.

\subsection{Colour Harmony}

Colour harmony concerns perceived compatibility among multiple colours in a composition. Classical harmony schemes (analogous, complementary, triadic, split-complementary, monochromatic) offer practical heuristics, while contemporary studies emphasize context dependence and user preference heterogeneity. In floral gift design, harmony is not only about formal balance but also about situational appropriateness; therefore, this thesis treats harmony as both a compositional property and a context-sensitive evaluation dimension.

\section{Emotion}

\subsection{Colour Imagery in Context}

Colour-evoked imagery is shaped by context, culture, and task. The same hue can signal different meanings across gifting scenarios (e.g., birthday, gratitude, apology, condolence), and these meanings can moderate emotional responses. Rather than assuming universal colour-emotion mappings, recent literature recommends contextualized interpretation and scenario-based testing.

\subsection{Emotion Alignment}

Emotion alignment refers to the degree to which a design's affective tone matches the intended message and usage context. In this thesis, alignment functions as a key mediating mechanism: when colour composition aligns with intended emotion, users are more likely to perceive the design as appropriate, attractive, and trustworthy. Misalignment, by contrast, may increase ambiguity and reduce evaluative confidence.

\section{Intentions}

\subsection{Design Attractiveness}

Design attractiveness captures immediate evaluative appeal and is frequently linked to approach-oriented responses. In human-computer interaction and consumer research, attractive visuals can increase willingness to explore, perceived usability, and overall satisfaction. For floral gifts, attractiveness is expected to be a proximal outcome of successful colour-emotion design.

\subsection{Adoption Intentions}

Adoption intention represents users' willingness to select, recommend, or use a generated design. Prior studies indicate that adoption intentions are jointly determined by aesthetic quality, emotional resonance, perceived fit to purpose, and trust in the design process. This thesis models adoption intention as a downstream variable influenced by colour composition and emotional alignment, directly and indirectly through attractiveness.

\section{Research Design}

\subsection{Floral Gift Design as Research Context}

Floral gift design provides an ecologically valid context where colour decisions are central and emotionally consequential. Compared with abstract palettes, floral compositions allow participants to evaluate colour in relation to recognizable forms, occasions, and social intentions.

\subsection{Quantitative Measures and Task Design}

The quantitative component uses structured rating tasks to measure perceived harmony, emotional alignment, attractiveness, and adoption intention. Stimuli are presented under controlled conditions, and measurement items are adapted from prior design, UX, and consumer-intention scales. This setup enables hypothesis testing on the relationships among colour features, emotion, and intention outcomes.

To keep measurement logic transparent, \autoref{tab:review-constructs} summarizes the core constructs, operational definitions, and expected roles in the model.

\begin{table}[htbp]
  \centering
  \begin{threeparttable}
    \caption{Core constructs in the Chapter 4 literature synthesis.}
    \label{tab:review-constructs}
    \begin{tabularx}{0.92\textwidth}{>{\raggedright\arraybackslash}p{3.1cm}>{\raggedright\arraybackslash}p{4.2cm}>{\raggedright\arraybackslash}X}
      \toprule
      \textbf{Construct} & \textbf{Operational Focus} & \textbf{Expected Role} \\
      \midrule
      Colour composition & Primary--secondary--accent structure; harmony and contrast control & Independent design factor that shapes perception and emotional tone. \\
      Emotion alignment & Match between intended gifting meaning and perceived affect & Mediating mechanism linking visual design and behavioral response. \\
      Design attractiveness & Immediate aesthetic appeal and perceived quality & Proximal outcome and partial mediator of intention. \\
      Adoption intention & Willingness to choose, recommend, or use a generated design & Downstream dependent variable reflecting practical acceptance. \\
      \bottomrule
    \end{tabularx}
  \end{threeparttable}
\end{table}

\subsection{Colour Representation \& Measurement}

Colour conditions are operationalized with explicit parameter controls and recorded metadata, enabling reproducibility across generated designs. Attribute-level measurements (hue, saturation/chroma, lightness/value, and contrast distribution) are paired with subjective ratings to connect computational descriptors and human perception.

\subsection{Qualitative and Mixed-Methods for Contextual Emotion and Usability}

Because emotional interpretation is context-sensitive, qualitative feedback (open-ended responses and brief interviews) complements quantitative scores. A mixed-methods strategy helps explain why specific colour compositions succeed or fail in different gifting scenarios and provides richer guidance for design iteration.

\subsection{Generative Systems: Controllability \& Bias}

Recent generative systems expand creative exploration but introduce concerns around controllability, consistency, and embedded aesthetic bias. This thesis therefore documents prompt and parameter settings, evaluates controllability at the attribute level, and reflects on possible dataset- or model-induced bias in colour outcomes.
